\documentclass[11pt]{article}
\usepackage[margin=1in]{geometry}
\usepackage{amsmath,amssymb,amsthm}
\usepackage{booktabs}
\usepackage{listings}
\usepackage{xcolor}
\usepackage[numbers]{natbib}
\usepackage{hyperref}

\newtheorem{theorem}{Theorem}
\newtheorem{lemma}[theorem]{Lemma}
\theoremstyle{definition}
\newtheorem{definition}[theorem]{Definition}
\theoremstyle{remark}
\newtheorem{remark}[theorem]{Remark}

\lstdefinelanguage{lean}{
  morekeywords={theorem,lemma,def,abbrev,namespace,end,open,import,by,fun,
    have,show,intro,exact,apply,rw,simp,simp_all,unfold,rintro,obtain,
    constructor,left,right,subst,rcases,by_cases,set,with,at,only,ring,
    if,then,else,rfl},
  sensitive=true,
  morecomment=[l]{--},
  morecomment=[s]{/-}{-/},
}
\lstset{
  language=lean,
  basicstyle=\ttfamily\footnotesize,
  keywordstyle=\color{blue!60!black}\bfseries,
  commentstyle=\color{green!40!black}\itshape,
  columns=fullflexible,
  breaklines=true,
  frame=single,
  extendedchars=true,
  literate=
    {∈}{{$\in$}}1 {ℤ}{{$\mathbb{Z}$}}1 {ℕ}{{$\mathbb{N}$}}1
    {·}{{$\cdot$}}1 {×}{{$\times$}}1 {∑}{{$\sum$}}1
    {⟨}{{$\langle$}}1 {⟩}{{$\rangle$}}1 {∀}{{$\forall$}}1
    {∨}{{$\vee$}}1 {←}{{$\leftarrow$}}1 {→}{{$\rightarrow$}}1
    {↔}{{$\leftrightarrow$}}1
}

\title{A Lean Proof of the Bit-Quad Euler Formula for 8-Connected Rasters:\\
the First Rung of a Raster Bridge for \emph{Trees Cannot Cut}}
\author{Justin Hart\\Viridis LLC}
\date{October 2026}

\begin{document}
\maketitle

\begin{abstract}
\emph{Trees Cannot Cut} argues that fuel breaks and habitat corridors split a
landscape into compartments only through their cycles. Its machine-checked core
carries the topological step (compartments equal circuit rank plus one) as an
explicit hypothesis. Project P1 aims to replace that hypothesis, on finite
rasters, with a proved identity: enclosed non-barrier 4-components equal
$C_8(B) - \chi(K B)$, where $K B$ is the cubical complex of the barrier raster
$B$. This paper does the first rung of that ladder. We give a Lean~4/Mathlib
proof of Gray's bit-quad formula in its 8-connected form: for every finite
raster $B \subset \mathbb{Z}^2$,
$4\chi(K B) = Q_1 - Q_3 - 2Q_D$, where $Q_1$, $Q_3$ and $Q_D$ count the
$2\times2$ windows holding exactly one pixel, exactly three pixels, or exactly
two diagonally opposite pixels of $B$. The proof double-counts over the vertices
of $K B$. Shift bijections show that $4F$, $2E$ and $V$ are sums of local window
quantities. A 16-case analysis then settles the pointwise identity
$4 - 2\,\mathrm{arms}(v) + n(v) = [n=1] - [n=3] - 2[\mathrm{diag}]$, which holds on
every non-empty window (it fails on the empty window, which never sits at a
vertex). The development is about 160 lines, uses no custom axioms, and
type-checks against a sealed statement module.
We explain how L1 feeds the single-pixel insertion lemma (L2) and the bridge (L3).
\end{abstract}

\section{Introduction}

Hart's \emph{Trees Cannot Cut}~\citep{Hart2026} makes a structural claim about
managed forest landscapes. A corridor or fuel-break network separates the fuel
matrix into compartments only through its cycles, so an acyclic network
(a ``tree'') cannot cut. The manuscript's Lean core proves the graph-theoretic
consequences of this claim. The topological input is assumed as a hypothesis,
\texttt{FaceCount}: the number of compartments is circuit rank plus one. It is
justified informally by Alexander duality, which Mathlib does not have.

Fire and fuel data usually come as rasters, and on rasters the topology is
finite and combinatorial. Project P1 therefore aims to prove the following
identity for every finite barrier raster $B$, read with 8-connected barriers
and a 4-connected complement:
\[
  \#\{\text{enclosed 4-components of } \mathbb{Z}^2\setminus B\}
  \;=\; C_8(B) - \chi(K B).
\]
The plan is a three-rung ladder. L1 (this paper) is a local formula for
$\chi(K B)$. L2 shows that adding one pixel preserves
$\text{enclosed} - C_8 + \chi$. L3 obtains the bridge by induction on $|B|$.
L1 is what makes the change in $\chi$ under a single insertion computable from
the four windows around the new pixel.

\section{Background and related work}

Gray~\citep{Gray1971} introduced \emph{bit-quads}: the $2\times2$ windows of a
binary image. He showed that area, perimeter and the Euler number are
``locally countable'' from bit-quad statistics. For 8-connected foreground the
Euler number is $(Q_1 - Q_3 - 2Q_D)/4$. Reading the foreground as 4-connected
flips the sign of the diagonal term: two pixels that touch only at a corner then
form two components, and $(6 - 0 + 2)/4 = 2$. In the 8-connected reading the two pixels meet at a corner of the closed
squares, which is exactly the cubical complex $K B$ used here. The (8,4)
pairing of foreground and background connectivity is standard in digital
topology, where it is the pairing under which a digital Jordan curve property
holds; see the survey of Kong and Rosenfeld~\citep{KongRosenfeld1989}.

The formal development uses Lean~4~\citep{deMouraUllrich2021} and
Mathlib~\citep{Mathlib2020}. We know of no prior machine-checked proof of the
bit-quad formula, but we have not searched exhaustively, and we make no claim
of novelty.

\section{The sealed statement}

A pixel $(i,j)\in\mathbb{Z}^2$ denotes the closed square
$[i,i+1]\times[j,j+1]$. For a finite set $B$ of pixels the sealed module
\texttt{ViridisStatements.P1.BitQuad} defines:
\begin{itemize}
  \item \texttt{vertices B}: the union of the four corners of every pixel;
  \item \texttt{hEdges B}: horizontal unit edges, named by their left end. Pixel
        $(i,j)$ contributes $(i,j)$ and $(i,j+1)$;
  \item \texttt{vEdges B}: vertical unit edges, named by their lower end. Pixel
        $(i,j)$ contributes $(i,j)$ and $(i+1,j)$;
  \item $\chi = |V| - (|E_h| + |E_v|) + |B|$, computed in $\mathbb{Z}$;
  \item the window at lattice point $v=(x,y)$, the bits
        $(a,b,c,d) = ([(x{-}1,y{-}1)\in B], [(x,y{-}1)\in B], [(x{-}1,y)\in B], [(x,y)\in B])$;
  \item $Q_1$, $Q_3$, $Q_D$: the number of $v\in$ \texttt{vertices B} whose window
        has one pixel, three pixels, or is a diagonal pair.
\end{itemize}
Every window that meets $B$ is centred on a vertex, so counting over
\texttt{vertices B} counts all non-empty windows.

\begin{theorem}[L1, \texttt{P1.BitQuad}]\label{thm:main}
For every finite $B\subset\mathbb{Z}^2$,
$4\,\chi(K B) = Q_1(B) - Q_3(B) - 2\,Q_D(B)$.
\end{theorem}

\section{Proof architecture}

\subsection{Double counting by shifts}
The only counting tool we need is the following lemma.

\begin{lemma}[\texttt{card\_filter\_shift}]
Let $V,S$ be finite subsets of $\mathbb{Z}^2$, and let $i,j$ be mutually inverse
maps of $\mathbb{Z}^2$ with $j(S)\subseteq V$. Then
$|\{v\in V : i(v)\in S\}| = |S|$.
\end{lemma}
Indicator sums are filter cardinalities (Mathlib's \texttt{Finset.sum\_boole}),
so with $V=\texttt{vertices B}$ the lemma gives:
\begin{align*}
  4|B| &= \sum_{v\in V} \big([a]+[b]+[c]+[d]\big)
  && \text{(each pixel is a window entry at each of its 4 corners)},\\
  2|E_h| &= \sum_{v\in V} \big([v\in E_h] + [(x{-}1,y)\in E_h]\big)
  && \text{(right and left arms)},\\
  2|E_v| &= \sum_{v\in V} \big([v\in E_v] + [(x,y{-}1)\in E_v]\big)
  && \text{(up and down arms)},\\
  |V| &= \sum_{v\in V} 1.
\end{align*}
Each line uses four (or two) shift bijections. The side condition
$j(S)\subseteq V$ says that every pixel's corners, and both endpoints of every
edge, are vertices. Each of these is a small membership lemma.

\subsection{Arms in terms of window bits}
Unfolding the edge sets gives $e\in E_h \iff e\in B \lor (e_1,e_2-1)\in B$ and
$e\in E_v \iff e\in B \lor (e_1-1,e_2)\in B$. At $v$ the four arms are therefore
$d\lor b$ (right), $c\lor a$ (left), $d\lor c$ (up) and $b\lor a$ (down).

\subsection{The local identity}
Summing the displays,
$4\chi = \sum_{v\in V} \big(4 - 2\,\mathrm{arms}(v) + n(v)\big)$ with
$n = [a]+[b]+[c]+[d]$. So it suffices to show, for each vertex,
\[
  4 - 2\,\mathrm{arms}(v) + n(v) = [n=1] - [n=3] - 2[\mathrm{diag}].
\]
Table~\ref{tab:local} lists the pattern classes. The identity \emph{fails} on
the empty window ($4\neq0$). We therefore use a lemma
(\texttt{window\_nonempty}): each vertex is a corner of some pixel of $B$, so
its window is non-empty. In Lean we split on the four membership propositions
(16 cases) and close each case with \texttt{simp\_all}. The brief suggested
\texttt{decide} on Booleans. Our first draft took that route, and its
\texttt{decide} correctly reported the unrestricted identity false. That is how
we found the empty-window case.

\begin{table}[h]
\centering
\begin{tabular}{lccccc}
\toprule
Pattern class & $n$ & arms & $4-2\,\mathrm{arms}+n$ & $[n{=}1]-[n{=}3]-2[\mathrm{diag}]$\\
\midrule
empty (never at a vertex) & 0 & 0 & 4 & 0\\
one pixel & 1 & 2 & 1 & 1\\
adjacent pair & 2 & 3 & 0 & 0\\
diagonal pair & 2 & 4 & $-2$ & $-2$\\
three pixels & 3 & 4 & $-1$ & $-1$\\
full & 4 & 4 & 0 & 0\\
\bottomrule
\end{tabular}
\caption{Local summand by window pattern class.}
\label{tab:local}
\end{table}

\section{Results: the Lean development}

The file \texttt{proofs/P1.BitQuad.lean} proves
\texttt{Candidate.main : ViridisStatements.P1.BitQuad.Statement}. It has 159
lines and imports only Mathlib and the sealed module. The project's local
checker (denylist, compile, statement match, axioms, kernel replay) reported
that it passed, with axioms \texttt{propext}, \texttt{Classical.choice} and
\texttt{Quot.sound}. The authoritative judgement is the verifier's
certification page. No Aristotle job was needed. The key lemmas follow.

\begin{lstlisting}
theorem card_filter_shift (V S : Finset (ℤ × ℤ)) (i j : ℤ × ℤ → ℤ × ℤ)
    (hij : ∀ p, i (j p) = p) (hji : ∀ v, j (i v) = v)
    (hS : ∀ p ∈ S, j p ∈ V) :
    (V.filter fun v => i v ∈ S).card = S.card := by
  apply Finset.card_nbij' i j
  · intro v hv
    simp only [Finset.coe_filter, Set.mem_setOf_eq] at hv
    exact hv.2
  · intro p hp
    simp only [Finset.coe_filter, Set.mem_setOf_eq]
    exact ⟨hS p hp, by rw [hij]; exact hp⟩
  · intro v _; exact hji v
  · intro p _; exact hij p

theorem mem_hEdges (B : Finset Pixel) (e : ℤ × ℤ) :
    e ∈ hEdges B ↔ e ∈ B ∨ (e.1, e.2 - 1) ∈ B

theorem window_nonempty (B : Finset Pixel) (v : ℤ × ℤ) (hv : v ∈ vertices B) :
    (v.1 - 1, v.2 - 1) ∈ B ∨ (v.1, v.2 - 1) ∈ B ∨ (v.1 - 1, v.2) ∈ B ∨ (v.1, v.2) ∈ B
\end{lstlisting}
The end of the main proof turns the global identity into the pointwise one:
\begin{lstlisting}
  rw [key, hQ1, hQ3, hQD, Finset.mul_sum, ← Finset.sum_sub_distrib,
      ← Finset.sum_sub_distrib]
  apply Finset.sum_congr rfl
  intro v hv
  have hne := window_nonempty B v hv
  simp only [mem_hEdges, mem_vEdges]
  unfold windowCount isDiagonal bits
  by_cases h1 : (v.1 - 1, v.2 - 1) ∈ B <;> by_cases h2 : (v.1, v.2 - 1) ∈ B <;>
    by_cases h3 : (v.1 - 1, v.2) ∈ B <;> by_cases h4 : (v.1, v.2) ∈ B <;>
    simp_all
\end{lstlisting}
Here \texttt{key} rewrites $4\,\chi$ as the sum of the local summand over
\texttt{vertices B}. It is obtained by distributing the sum and substituting
the eight shift identities. The proof concludes with \texttt{ring}.

\section{How L1 serves the bridge}

L2 asks that inserting a pixel $p\notin B$ preserve
$\text{enclosed} - C_8 + \chi$. By Theorem~\ref{thm:main}, $\Delta\chi$ is
determined by the four windows centred on the corners of $p$. Only those
windows change, so
$4\,\Delta\chi = \Delta Q_1 - \Delta Q_3 - 2\,\Delta Q_D$ is a finite local
computation over $2^{8}$ neighbourhood patterns of $p$. The non-local part of
L2 is the change in $C_8$ and in the number of enclosed components. Two
8-components merging through $p$ can close off a new region. That part is not
touched here. We expect that the natural Lean statement of L2 will be phrased
in terms of the local $\Delta\chi$ that L1 supplies. The three controls listed
in the project README (the 3$\times$3 ring witness, the 2$\times$2 block and
the four-arm plus) give concrete checks: the sealed module verifies
$\chi=1$ for the 2$\times$2 block and $\chi=0$ for the ring by computation.

\section{Limitations}

\begin{itemize}
  \item L1 is a statement about the cubical complex counted combinatorially. It
        says nothing about connected components or enclosure. The topological
        content of the bridge sits entirely in L2.
  \item The statement fixes the 8-connected reading of the foreground. The
        4-connected variant ($+2Q_D$) is not formalised here.
  \item L1 by itself does not discharge the \texttt{FaceCount} hypothesis of
        \emph{Trees Cannot Cut}. It also does not address the raster-to-graph
        correspondence, whose gap the 2$\times$2 block control exposes
        (circuit rank 3, nothing enclosed).
\end{itemize}

\section{Reproducibility}

Toolchain \texttt{leanprover/lean4:v4.28.0}; Mathlib commit
\texttt{8f9d9cff6bd728b17a24e163c9402775d9e6a365}. The sealed statement is
\texttt{P1.BitQuad.lean} (SHA-256 \texttt{cd3cdeba48f67f95\allowbreak45fb6945c87037893aa98b018922d3d0b7220cc389c98e23}).
To reproduce, run
\texttt{lean proofs/P1.BitQuad.lean} with \texttt{LEAN\_PATH} including Mathlib
and the sealed modules. Compilation takes about 10~s on the session machine. The
proof contains no \texttt{sorry}, \texttt{native\_decide} or custom axioms.

\section{Changelog}
\begin{itemize}
  \item 2026-10-27 (P1-T001): first version. Lean proof of L1 (\texttt{P1.BitQuad});
        proof architecture; role of L1 in the L2/L3 ladder.
\end{itemize}

\bibliographystyle{plainnat}
\bibliography{refs}
\end{document}
